
% Authoritative LaTeX manuscript; Markdown is a generated compatibility projection.

% <!-- part-title: The Profession -->

% <!-- chapter-title: What Cannot Be Delegated? -->

% <!-- coda: true -->

\chapter{What Cannot Be Delegated?}

Making implementation cheaper does not make software engineering easy; it changes where the difficulty lies.

This chapter has followed that change from software work to engineering traditions, beyond software to other consequential professions, and finally to the engineer’s own expertise and authority. Across these settings, the same distinction persists. Work can be delegated, and authority to make particular decisions can be granted to people or mechanisms. Responsibility for the resulting engineering arrangement remains with those who authorize and govern it.

That distinction does not reserve a permanent set of tasks for humans. Better machines may eventually perform nearly every activity discussed in this book, including constructing models, evaluating evidence, proposing controls, and exercising consequential judgment. What matters is whether people and institutions retain sufficient understanding, control, and authority to remain answerable for the resulting systems.

The Preface asked what we can safely delegate. The Conclusion approaches the same boundary from the other direction:

What cannot be delegated?
